\documentclass{article}
\usepackage[margin=1.5cm]{geometry}
\usepackage{amsmath}
\usepackage{enumitem}

\setlist{nosep,leftmargin=*}

\begin{document}
\twocolumn

\title{Chapter 9 Warm-Up: Organizing Course Data with Lists}
\author{Prof. Jordan C. Hanson}
\date{}

\maketitle

\section{Nested Course Lists}

\begin{enumerate}

\item Chapter 9 develops techniques for modifying, sorting, and searching lists.  We will use a nested list to represent a college course schedule. One course is represented by

\begin{verbatim}
course = ["PHYS",100,"SLC", 202]
\end{verbatim}

The four elements represent

\begin{verbatim}
[subject, number, building, room]
\end{verbatim}

Thus, this record represents Physics 100 in SLC 202.

\end{enumerate}


\section{A List of Lists}

\begin{enumerate}

\item Several course records can be stored inside one master list:

\begin{verbatim}
courses = [
 ["PHYS", 100, "SLC", 202],
 ["MATH", 150, "SLC", 105],
 ["BIOL", 101, "SLC", 210],
 ["PHYS", 200, "SLC", 202]
]
\end{verbatim}

Each element of \texttt{courses} is itself a list.  Print one complete record:

\begin{verbatim}
print(courses[1])
\end{verbatim}

Then print only its room number:

\begin{verbatim}
print(courses[1][3])
\end{verbatim}

\end{enumerate}


\section{Appending Courses}

\begin{enumerate}

\item Define a function that adds a course to the master list:

\begin{verbatim}
def add_course(courses, course):
    courses.append(course)
\end{verbatim}

Test it using

\begin{verbatim}
new_course = ["COSC",120,"SLC",416]
add_course(courses, new_course)
\end{verbatim}

Use \verb+print(len(courses))+ to verify that the new record was added.

\end{enumerate}


\section{Sorting the Master List}

\begin{enumerate}

\item Lists can be sorted using \texttt{sort()}. Define

\begin{verbatim}
def sort_courses(courses):
    courses.sort()
\end{verbatim}

Then run

\begin{verbatim}
sort_courses(courses)
for course in courses:
    print(course)
\end{verbatim}

Since the subject is element \texttt{0}, Python first sorts the records alphabetically by subject. If two subjects are the same, which element is compared next?

\end{enumerate}


\section{Searching for Room Conflicts}

\begin{enumerate}

\item Two courses conflict if they use the same building and room. Define

\begin{verbatim}
def same_room(c1, c2):
    return c1[2] == c2[2] and c1[3] == c2[3]
\end{verbatim}

Now use nested loops to compare every pair:

\begin{verbatim}
for i in range(len(courses)):
    for j in range(i+1,len(courses)):
        if same_room(courses[i],courses[j]):
            print("Room conflict:")
            print(courses[i])
            print(courses[j])
\end{verbatim}

Why does the inner loop begin at \texttt{i+1} instead of zero?

\end{enumerate}

\section{Warm-Up Exercise}

\begin{enumerate}

\item Add the following records using
\texttt{add\_course()}:

\begin{verbatim}
["MATH", 250, "SLC", 105]
["CHEM", 110, "SLC", 210]
["PHYS", 150, "SLC", 310]
\end{verbatim}

Complete the following:

\begin{itemize}
\item Sort the master list
\item Print every course with a \texttt{for} loop
\item Use the nested-loop search to print every room conflict
\end{itemize}

Before running each bit of code, predict which new courses will cause room conflicts.

\end{enumerate}

\end{document}